\chapter{AHP Pairwise Comparison Matrix and Step-by-Step Calculation}\label{app:ahp-matrices}
% Nguon: 03_phan-tich/ahp/2026-07-11_ahp-matran-final_v1.0_final_user.csv
% Mot ma tran duy nhat (1 chuyen gia dai dien) — khong co tong hop nhieu chuyen gia.
% Tat ca cac buoc duoi day tai lap chinh xac thuat toan power-iteration cua
% .claude/skills/mcdm-toolkit/scripts/ahp.py (eigen_weights + consistency), khong tinh tay.

\Cref{tab:pcm-full} reproduces the full $7 \times 7$ pairwise comparison matrix elicited
from the representative expert as described in \Cref{sec:stage2}. Entries above the
diagonal are the 21 elicited judgements on the Saaty scale; entries below the diagonal
follow from reciprocity.

\begin{table}[htbp]
  \centering
  \caption{Pairwise comparison matrix of the seven retained criteria.}
  \label{tab:pcm-full}
  \begin{tabular}{@{}lccccccc@{}}
    \toprule
     & \textbf{Price} & \textbf{CPU} & \textbf{RAM} & \textbf{SSD} & \textbf{Brand} & \textbf{Size/weight} & \textbf{Battery} \\
    \midrule
    \textbf{Price}       & 1    & 1    & 5    & 4    & 7    & 8    & 8 \\
    \textbf{CPU}         & 1    & 1    & 3    & 2    & 7    & 9    & 8 \\
    \textbf{RAM}         & 1/5  & 1/3  & 1    & 1/2  & 5    & 6    & 6 \\
    \textbf{SSD}         & 1/4  & 1/2  & 2    & 1    & 5    & 7    & 6 \\
    \textbf{Brand}       & 1/7  & 1/7  & 1/5  & 1/5  & 1    & 2    & 3 \\
    \textbf{Size/weight} & 1/8  & 1/9  & 1/6  & 1/7  & 1/2  & 1    & 1 \\
    \textbf{Battery}     & 1/8  & 1/8  & 1/6  & 1/6  & 1/3  & 1    & 1 \\
    \bottomrule
  \end{tabular}
  \floatnote{\emph{Note:} entries above the diagonal are the 21 elicited judgements on the Saaty scale; entries below the diagonal follow from reciprocity.}{\emph{Source:} Authors.}
\end{table}

\section{Step 1 --- Column Sums}\label{app:ahp-step1}

Each column of the matrix $A$ in \Cref{tab:pcm-full} is summed to obtain the normalising
constant used in the following step. \Cref{tab:ahp-colsum} reports the seven column
totals.

\begin{table}[htbp]
  \centering
  \caption{Column sums of the pairwise comparison matrix.}
  \label{tab:ahp-colsum}
  \begin{tabular}{@{}lccccccc@{}}
    \toprule
     & \textbf{Price} & \textbf{CPU} & \textbf{RAM} & \textbf{SSD} & \textbf{Brand} & \textbf{Size/weight} & \textbf{Battery} \\
    \midrule
    Column sum & 2.8429 & 3.2123 & 11.5333 & 8.0095 & 25.8333 & 34.0000 & 33.0000 \\
    \bottomrule
  \end{tabular}
  \floatsource{\emph{Source:} Authors.}
\end{table}

\section{Step 2 --- Weight Vector by Power Iteration (Eigenvector Method)}\label{app:ahp-step2}

Following the justification given in \Cref{sec:stage2} for using the principal eigenvector
rather than the normalised column-average approximation
\parencite{saaty_decisionmaking_2003,ishizaka_how_2006}, the weight vector $\mathbf{w}$ is
obtained by power iteration: starting from the uniform vector
$\mathbf{w}^{(0)} = (1/7, \dots, 1/7)$, each iteration computes
$\mathbf{w}^{(k+1)} = A\mathbf{w}^{(k)} / \left\lVert A\mathbf{w}^{(k)} \right\rVert_1$
until the vector stabilises. \Cref{tab:ahp-iterations} traces the first six iterations;
the vector has already converged to four decimal places by the sixth iteration, which is
taken as the final weight vector.

\begin{table}[htbp]
  \centering
  \caption{Convergence of the power-iteration weight vector across six iterations.}
  \label{tab:ahp-iterations}
  \begin{tabular}{@{}lccccccc@{}}
    \toprule
    \textbf{Iteration} & \textbf{Price} & \textbf{CPU} & \textbf{RAM} & \textbf{SSD} & \textbf{Brand} & \textbf{Size/weight} & \textbf{Battery} \\
    \midrule
    $k=1$ & 0.2871 & 0.2618 & 0.1607 & 0.1837 & 0.0565 & 0.0257 & 0.0246 \\
    $k=2$ & 0.3550 & 0.2734 & 0.1208 & 0.1621 & 0.0405 & 0.0242 & 0.0240 \\
    $k=3$ & 0.3522 & 0.2774 & 0.1182 & 0.1581 & 0.0424 & 0.0258 & 0.0259 \\
    $k=4$ & 0.3488 & 0.2771 & 0.1201 & 0.1591 & 0.0431 & 0.0259 & 0.0259 \\
    $k=5$ & 0.3490 & 0.2769 & 0.1202 & 0.1593 & 0.0430 & 0.0258 & 0.0258 \\
    $k=6$ & 0.3491 & 0.2769 & 0.1201 & 0.1593 & 0.0430 & 0.0258 & 0.0258 \\
    \bottomrule
  \end{tabular}
  \floatsource{\emph{Source:} Authors.}
\end{table}

The converged weight vector is therefore
\[
\mathbf{w} =
\bigl(0.3491,\ 0.2769,\ 0.1201,\ 0.1592,\
      0.0430,\ 0.0258,\ 0.0258\bigr),
\]
matching the values reported in \Cref{tab:ahp-weights} of \Cref{sec:weights}.

\section{Step 3 --- Consistency Check}\label{app:ahp-step3}

The converged weight vector is substituted back into $A\mathbf{w}$, and each component is
divided by the corresponding weight to obtain seven estimates of the principal eigenvalue,
as shown in \Cref{tab:ahp-consistency}. Because $A$ satisfies exact reciprocity and the
power iteration has converged, all seven ratios agree, giving
$\lambda_{\max} = 7.3823$ directly rather than as an average over inconsistent estimates.

\begin{table}[htbp]
  \centering
  \caption{Consistency check: $A\mathbf{w}$ and the ratio $(A\mathbf{w})_i / w_i$.}
  \label{tab:ahp-consistency}
  \begin{tabular}{@{}lccccccc@{}}
    \toprule
     & \textbf{Price} & \textbf{CPU} & \textbf{RAM} & \textbf{SSD} & \textbf{Brand} & \textbf{Size/weight} & \textbf{Battery} \\
    \midrule
    $(A\mathbf{w})_i$        & 2.5773 & 2.0444 & 0.8865 & 1.1756 & 0.3174 & 0.1903 & 0.1908 \\
    $(A\mathbf{w})_i / w_i$  & 7.3823 & 7.3823 & 7.3823 & 7.3823 & 7.3823 & 7.3823 & 7.3823 \\
    \bottomrule
  \end{tabular}
  \floatsource{\emph{Source:} Authors.}
\end{table}

With $\lambda_{\max} = 7.3823$ and $n = 7$, the consistency index from \Cref{eq:cr} is

\[
  CI = \frac{\lambda_{\max} - n}{n - 1} = \frac{7.3823 - 7}{6} = 0.0637 .
\]

Using the random index $RI = 1.32$ tabulated in \Cref{tab:ri} for $n = 7$, the consistency
ratio is

\[
  CR = \frac{CI}{RI} = \frac{0.0637}{1.32} = 0.0483 ,
\]

which is below the acceptance threshold of $0.10$ \parencite{Saaty2008DECISIONMW}, so the
judgement matrix is accepted without re-elicitation, as reported in \Cref{sec:weights}.

\section{Robustness Check --- Geometric-Mean Weights}\label{app:ahp-step4}

As a cross-check against the derivation method rather than against the judgement data
itself, the same matrix $A$ is also solved by the geometric-mean-of-rows method: for each
row $i$, $\tilde{w}_i = \left(\prod_{j=1}^{n} a_{ij}\right)^{1/n}$, and the seven values are
then normalised to sum to one. \Cref{tab:ahp-geomean} compares the two methods.

\begin{table}[htbp]
  \centering
  \caption{Eigenvector weights versus geometric-mean weights.}
  \label{tab:ahp-geomean}
  \begin{tabular}{@{}lccccccc@{}}
    \toprule
     & \textbf{Price} & \textbf{CPU} & \textbf{RAM} & \textbf{SSD} & \textbf{Brand} & \textbf{Size/weight} & \textbf{Battery} \\
    \midrule
    Eigenvector (adopted)  & 0.3491 & 0.2769 & 0.1201 & 0.1592 & 0.0430 & 0.0258 & 0.0258 \\
    Geometric mean         & 0.3366 & 0.2882 & 0.1185 & 0.1615 & 0.0429 & 0.0264 & 0.0259 \\
    \bottomrule
  \end{tabular}
  \floatsource{\emph{Source:} Authors.}
\end{table}

The two methods rank all seven criteria identically and differ from one another by at most
0.013 in absolute weight, confirming that the priority structure reported in
\Cref{sec:weights} does not depend on which of the two derivation methods is used.
